\addcontentsline{toc}{chapter}{CONCLUSIONS}
\chapter*{CONCLUSIONS}

In summary, the evidence presented supports the idea that a student-focused forum is a good idea for implementation, as it allows students to centralize and give them a convenient way to connect with each other and discuss topics of concern to them.

It's also worth noting that this platform requires a high level of security, as it uses students' personal information and their anonymous posts and comments. Various security measures were employed to ensure user safety. Specifically, the OAuth protocol was implemented, multi-factor authentication (MFA) support was implemented, and the Supabase platform, with its encryption and access control support, was used for data storage and processing.

On the frontend, Next.js was used to implement protection mechanisms against common attacks (XSS, CSRF), and secure cookie attributes were used. All API requests are validated, preventing the possibility of SQL injections and other vulnerabilities.

As a result, the project demonstrates the practical application of Secure Application Development principles, combining forum functionality with modern security measures, ensuring resilience to cyberthreats and meeting the stated requirements.

In addition to providing security and essential functionality, the project establishes the groundwork for future growth and expansion.  Other features that might be added to the forum include role-based moderating tools, interaction with university services, and sophisticated analytics to track user behavior and community involvement.  This demonstrates the project's usefulness as it has the potential to develop into a strong platform that fosters peer connections as well as academic cooperation and knowledge exchange inside the university.